% \iffalse meta-comment
%
% docxlike-table.dtx 是表格
%
% \fi
%
% \section{表格}
%
% \changes{v3.2a}{2026/09/28}{表格属性“表格→选项→默认单元格边距”（\texttt{table-properties/table/options/default-cell-margins}），
%   \pkg{tabularray} 的四个边距照它设；\pkg{tabularray} 起表与每页开的那几段不发撑杆}
% \changes{v3.2a}{2026/09/28}{\texttt{m} 列在整行的高里上下居中（\hologo{LuaLaTeX}，Word 的 \texttt{vAlign!="center"}）；
%   \cs{multirow} 不写位置时顶端对齐（Word 合并的格）；\pkg{tabularray} 的单根竖线不占宽、骑在分界上（探针 tbl1）}
% \changes{v3.2a}{2026/09/27}{表格属性“行”的“指定高度”（\texttt{table-properties/row/specify-height}）：
%   行高最小值，内容高了照内容，多出来的空在行的下面}
% \changes{v3.2a}{2026/09/26}{新增表格：\pkg{array} 一族（\env{tabular}、\env{tabularx}、
%   \env{longtable} 等）照 Word 的表格排：格边距 5.4\,bp，竖线不占宽，\pkg{booktabs} 线外不留空，
%   行的撑杆是 Word 的行盒，格里的字与正文一样吃网格增量，格里的段按“在表格单元格里”断行、
%   居中右对齐按整个格宽；\pkg{tabularray} 的表只补字宽}
%
% 配置单元格边距、行盒、边框和表内断行。
% \begin{itemize}
% \item 格左右各有 5.4\,bp 的边距（\texttt{Normal Table} 的 \texttt{tblCellMar} 是 108 缇），上下没有。
% \item 竖的框线画在格与格的分界上，不占宽，格里的字从分界起空一个格边距。横的框线占它自己的
%   粗细：一行是“上边框、各格的段”摞起来，最后一行下面再是下边框；五号字单倍贴网格的一行
%   是 19.55 加 0.5 的线，字的基线在线下 13.56。
% \item 格里的段照正文的规矩吃网格增量（汉字一份、西文半份、空格 0.25\,bp），居中与右对齐的
%   行宽也算上末字的增量；两端对齐的行撑满整个格宽，不按整数个网格装。
% \item 表整体往左让出一个格边距，格里的字与正文左缘对齐。
% \end{itemize}
%
% \hologo{LaTeX} 的表格宏包各有一套自己的量，这里把它们换成 Word 的：
% \begin{itemize}
% \item \cs{tabcolsep} 取 Word 的格边距 5.4\,bp（装包时设，文档里改了照改的）。
% \item \cs{vline}（\pkg{array} 的 \texttt{\textbar} 也是它）画在分界线上、左右各退一半线宽，不占宽。
%   横线（\cs{hline}、\cs{cline}、\pkg{booktabs} 的三条线）照旧占高，\pkg{booktabs} 线上线下的空
%   归零：Word 的框线外面没有这一截。
% \item 行的撑杆换成当前那一段的行盒（贴网格时一格），字在行里的高低与 Word 一样。
% \item 格里的段按“在表格单元格里”断行（\cs{l\_docxlike\_table\_cell\_bool}）；格里写的
%   \cs{centering}、\cs{raggedright}、\cs{raggedleft} 也算照样式设的对齐，交给断行层排，行宽
%   算上末字的增量。\texttt{l}、\texttt{c}、\texttt{r} 列的格是一个盒，没有段，断行层在装盒时
%   给格里的字补上增量与字偶距（属性 \cs{docxlike@cell@attr} 标着哪些字在格里）。
% \end{itemize}
%
% \pkg{tabularray} 只做后一条的字宽与竖线：它自己量每一格的盒来定行高与竖直居中，撑杆与格里的段
% 换成 Word 的行盒后，居中的格与别的格基线对不齐。所以它的表不换撑杆、格里的段照旧交回
% \hologo{TeX}（字宽照样补）；单根竖线改成不占宽（见本节末尾）。
%
% \texttt{m} 列在 \pkg{array} 里是对着行的基线居中，跟 Word 的“垂直居中”（对着整行居中）
% 不一样：同一行里有折成几行的 \texttt{p} 格时，\texttt{m} 格落在第一行上。\hologo{LuaLaTeX} 上
% 一行排完再挪到整行正中（见 \texttt{m} 列那一段），别的引擎照 \pkg{array}。
% 只在 Word 版面（\cs{g\_docxlike\_word\_layout\_bool}）下做；网格增量与断行只在
% \hologo{LuaLaTeX} 上有。
%    \begin{macrocode}
%<*docxlike-table>
\bool_new:N \l_docxlike_table_bool
\bool_set_true:N \l_docxlike_table_bool
\bool_new:N \l_docxlike_table_cell_bool
\dim_set:Nn \tabcolsep { 5.4bp }
\sys_if_engine_luatex:T { \newattribute \docxlike@cell@attr \newattribute \docxlike@mcell@attr }
\box_new:N \l__docxlike_table_strut_box
\bool_new:N \l__docxlike_table_tblr_bool
%    \end{macrocode}
%
% \cs{l\_docxlike\_table\_bool} 关掉就不换：文档类拿 \env{tabular} 当排版工具用的地方（堆几行字、
% 对齐封面上的字段），Word 里那不是表，局部关掉；整篇不照 Word 排的文档全局关掉。
%
% \begin{macro}{\docxlike_table_begin:}
% 表格环境开头调一次，设置都在环境的组里。文档类给自己的表另设字体时，在这之后设，撑杆
% 再按新的字体取一次（\cs{docxlike\_format\_strut\_box:N}）。\cs{l\_\_docxlike\_table\_tblr\_bool}
% 是 \pkg{tabularray} 的表，见上。
%    \begin{macrocode}
\cs_new_protected:Npn \docxlike_table_begin:
  {
    \bool_lazy_and:nnT { \g_docxlike_word_layout_bool } { \l_docxlike_table_bool }
      {
        \bool_lazy_and:nnT
          { ! \l__docxlike_table_tblr_bool }
          { \dim_compare_p:nNn { \l_docxlike_table_row_height_dim } > { \c_zero_dim } }
          { \cs_set_protected:Npn \@arstrut { \tex_unhcopy:D \@arstrutbox \__docxlike_table_row_min: } }
        \bool_set_true:N \l_docxlike_table_cell_bool
        \bool_if:NF \l__docxlike_table_tblr_bool
          { \bool_set_true:N \l_docxlike_linebreak_align_bool }
        \cs_set_protected:Npn \vline
          {
            \tex_hskip:D -0.5 \arrayrulewidth
            \tex_vrule:D width \arrayrulewidth
            \tex_hskip:D -0.5 \arrayrulewidth
          }
        \cs_if_exist:NT \aboverulesep
          {
            \dim_zero:N \aboverulesep
            \dim_zero:N \belowrulesep
            \dim_zero:N \abovetopsep
            \dim_zero:N \belowbottomsep
          }
        \bool_if:NF \l__docxlike_table_tblr_bool
          {
            \docxlike_format_strut_box:N \l__docxlike_table_strut_box
            \box_set_eq:NN \strutbox \l__docxlike_table_strut_box
            \cs_if_exist:NT \@multirow
              { \cs_set_protected:Npn \multirow { \peek_meaning:NTF [ { \@multirow } { \@multirow [ t ] } } }
            \sys_if_engine_luatex:T
              {
                \cs_if_exist:NT \ar@align@mcell
                  {
                    \cs_set_protected:Npn \ar@align@mcell
                      { \docxlike@mcell@attr = 1 \scan_stop: \hbox:n { \box_use_drop:N \ar@mcellbox } }
                  }
              }
          }
        \sys_if_engine_luatex:T { \docxlike@cell@attr = 1 \scan_stop: }
      }
  }
%    \end{macrocode}
% \end{macro}
%
% \subsection{行高}
%
% 表格属性“行”那一页的“指定高度”：\cs{docxlikesetup}\texttt{\{table-properties=\{row=\{specify-height=}\meta{长度}\texttt{\}\}\}}，
% 之后（同一组里）的表每一行至少这么高，内容高了照内容。Word 勾上“指定高度”时“行高值是”默认
% “最小值”，存成 \texttt{w:trHeight}，量的是两道横线之间的空（线宽另算）；“固定值”没有做。
% 字在行里的高低不变，多出来的空在行的下面（Word 格里默认顶端对齐）。
%
% 落实的办法：\pkg{array} 在每一行的头一格（\cs{multicolumn} 的格、\pkg{longtable} 也一样）先放
% \cs{@arstrut}，把它换成撑杆后面再跟一道零宽的竖线，深度补到行高减撑杆的高。行的高深取各格
% 的最大值，\texttt{p} 格折成几行时它自己的深度大于这道线，行高就照内容，与 Word 一样。
% \pkg{tabularray} 的表不做，它有自己的 \texttt{rows=\{ht\}}。
%    \begin{macrocode}
\dim_new:N \l_docxlike_table_row_height_dim
\keys_define:nn { docxlike }
  {
    table-properties .code:n =
      { \__docxlike_keys_set:nn { docxlike / table-properties } {#1} } ,
  }
\keys_define:nn { docxlike / table-properties }
  {
    row .code:n =
      { \__docxlike_keys_set:nn { docxlike / table-properties / row } {#1} } ,
    table .code:n =
      { \__docxlike_keys_set:nn { docxlike / table-properties / table } {#1} } ,
    unknown .code:n = { \__docxlike_unknown_key: } ,
  }
%    \end{macrocode}
%
% 表格属性“表格”那一页“选项”里的“默认单元格边距”：\texttt{table-properties=\{table=\{options=\{default-cell-margins=\{}%
% \texttt{left=\ldots, right=\ldots, top=\ldots, bottom=\ldots\}\}\}\}}。Word 的默认值是左右 5.4\,bp、上下 0
% （\texttt{Normal Table} 的 \texttt{tblCellMar}）。\pkg{array} 一族只有一个 \cs{tabcolsep}，左右一样，取
% \option{left}；上下的没有地方放。\pkg{tabularray} 的四个都有：进表时照这四个设它的 \texttt{leftsep}、
% \texttt{rightsep}、\texttt{abovesep}、\texttt{belowsep}（它默认左右 6\,pt、上下 2\,pt），表里自己写的照写的。
%    \begin{macrocode}
\dim_new:N \l_docxlike_table_margin_left_dim
\dim_new:N \l_docxlike_table_margin_right_dim
\dim_new:N \l_docxlike_table_margin_top_dim
\dim_new:N \l_docxlike_table_margin_bottom_dim
\dim_set:Nn \l_docxlike_table_margin_left_dim { 5.4bp }
\dim_set:Nn \l_docxlike_table_margin_right_dim { 5.4bp }
\keys_define:nn { docxlike / table-properties / table }
  {
    options .code:n =
      { \__docxlike_keys_set:nn { docxlike / table-properties / table / options } {#1} } ,
    unknown .code:n = { \__docxlike_unknown_key: } ,
  }
\keys_define:nn { docxlike / table-properties / table / options }
  {
    default-cell-margins .code:n =
      { \__docxlike_keys_set:nn { docxlike / table-properties / table / options / margins } {#1} } ,
    unknown .code:n = { \__docxlike_unknown_key: } ,
  }
\keys_define:nn { docxlike / table-properties / table / options / margins }
  {
    left .code:n =
      {
        \dim_set:Nn \l_docxlike_table_margin_left_dim { \__docxlike_bp:n {#1} }
        \dim_set_eq:NN \tabcolsep \l_docxlike_table_margin_left_dim
      } ,
    right .code:n = { \dim_set:Nn \l_docxlike_table_margin_right_dim { \__docxlike_bp:n {#1} } } ,
    top .code:n = { \dim_set:Nn \l_docxlike_table_margin_top_dim { \__docxlike_bp:n {#1} } } ,
    bottom .code:n = { \dim_set:Nn \l_docxlike_table_margin_bottom_dim { \__docxlike_bp:n {#1} } } ,
    unknown .code:n = { \__docxlike_unknown_key: } ,
  }
\keys_define:nn { docxlike / table-properties / row }
  {
    specify-height .dim_set:N = \l_docxlike_table_row_height_dim ,
    unknown .code:n = { \__docxlike_unknown_key: } ,
  }
\cs_new_protected:Npn \__docxlike_table_row_min:
  {
    \dim_compare:nNnT
      { \l_docxlike_table_row_height_dim } > { \box_ht:N \@arstrutbox }
      {
        \tex_vrule:D height \c_zero_dim width \c_zero_dim
          depth \dim_eval:n { \l_docxlike_table_row_height_dim - \box_ht:N \@arstrutbox }
      }
  }
%    \end{macrocode}
%
% \emph{\texttt{m} 列对着整行居中。} Word 的“垂直居中”（\texttt{vAlign="center"}）是在整行的高里居中（探针 tbl1 的 T4：
% “居中”落在同一行“上”“下”两行正中）；\pkg{array} 的 \texttt{m} 列只对着行的基线摆。\hologo{LuaLaTeX} 上
% 把 \texttt{m} 格的盒子打上 \cs{docxlike@mcell@attr}，一行排完（\texttt{hpack\_filter} 的 \texttt{fin\_row}）
% 由断行层按这一行各格的最高与最深把它挪到正中，见 \file{docxlike-linebreak.dtx}。别的引擎照 \pkg{array}。
%
% \emph{纵向合并的格顶端对齐。} Word 里合并起来的格与别的格一样默认顶端对齐（\texttt{vMerge} 的格没写
% \texttt{vAlign}，探针 tbl1 的 T4“合并”与同一行的“上”同一条基线），\pkg{multirow} 默认上下居中。
% 表里没写位置的 \cs{multirow} 改成 \texttt{t}，写了 \oarg{c} 的照写的（Word 的 \texttt{vAlign="center"}）。
%
% \pkg{array} 一族都在 \env{tabular} 或 \env{tabular*} 上搭，\pkg{tabularx}、\pkg{tabulary}
% 的环境先量后排，量的那一遍也要同样的字宽，所以各自的环境都挂。\env{array} 是公式里的，
% 不挂。\pkg{tabularray} 的表在装包之后挂，连同它 \texttt{booktabs} 库里的三个环境：那三个
% 也是 \cs{NewTblrEnviron} 定的表，漏挂的话格里的字照样补上字宽，量列宽的那一遍却没补，
% 没有空格可压的格就溢出。
%    \begin{macrocode}
\clist_map_inline:nn { tabular , tabular* , tabularx , tabulary , longtable , xltabular }
  { \hook_gput_code:nnn { env / #1 / begin } { docxlike-table } { \docxlike_table_begin: } }
\hook_gput_code:nnn { package / tabularray / after } { docxlike-table }
  {
    \clist_map_inline:nn { tblr , longtblr , talltblr , booktabs , longtabs , talltabs }
      {
        \hook_gput_code:nnn { env / #1 / begin } { docxlike-table }
          {
            \bool_set_true:N \l__docxlike_table_tblr_bool
            \docxlike_table_begin:
            \__docxlike_table_tblr_margins:n {#1}
          }
      }
    \__docxlike_table_tblr_vlines:
    \__docxlike_table_tblr_paras:
  }
\cs_new_protected:Npn \__docxlike_table_tblr_margins:n #1
  {
    \bool_lazy_and:nnT { \g_docxlike_word_layout_bool } { \l_docxlike_table_bool }
      {
        \SetTblrInner [#1]
          {
            leftsep = \dim_use:N \l_docxlike_table_margin_left_dim ,
            rightsep = \dim_use:N \l_docxlike_table_margin_right_dim ,
            abovesep = \dim_use:N \l_docxlike_table_margin_top_dim ,
            belowsep = \dim_use:N \l_docxlike_table_margin_bottom_dim ,
          }
      }
  }
%    \end{macrocode}
%
% tabularray 用于更新 pagetotal 和容纳整表的空段须禁用撑杆，否则会额外占据一行。
%
% 补丁打在 \pkg{tabularray} 的内部函数上，\cs{patchcmd} 按打补丁那一刻的 catcode 重扫宏体：下划线与冒号
% 设成字母，空格设回 10（宏体里有 \texttt{\string\hrule height \string~ 0pt}，\pkg{expl3} 语法下重扫会把那个空格
% 吞掉，比对就过不了），打完还原。打不上就提示一句。
%    \begin{macrocode}
\int_new:N \l__docxlike_table_catcode_us_int
\int_new:N \l__docxlike_table_catcode_colon_int
\int_new:N \l__docxlike_table_catcode_space_int
\msg_new:nnn { docxlike } { table-patch-failed }
  { Could~not~patch~\token_to_str:N #1~of~tabularray;~its~paragraphs~keep~their~struts. }
\cs_new_protected:Npn \__docxlike_table_patch:Nnn #1#2#3
  {
    \int_set:Nn \l__docxlike_table_catcode_us_int    { \char_value_catcode:n { 95 } }
    \int_set:Nn \l__docxlike_table_catcode_colon_int { \char_value_catcode:n { 58 } }
    \int_set:Nn \l__docxlike_table_catcode_space_int { \char_value_catcode:n { 32 } }
    \char_set_catcode_letter:n { 95 }
    \char_set_catcode_letter:n { 58 }
    \char_set_catcode_space:n  { 32 }
    \patchcmd #1 {#2} {#3} { }
      { \msg_warning:nnn { docxlike } { table-patch-failed } {#1} }
    \char_set_catcode:nn { 95 } { \l__docxlike_table_catcode_us_int }
    \char_set_catcode:nn { 58 } { \l__docxlike_table_catcode_colon_int }
    \char_set_catcode:nn { 32 } { \l__docxlike_table_catcode_space_int }
  }
\cs_new_protected:Npn \__docxlike_table_tblr_paras:
  {
    \bool_lazy_all:nT
      {
        { \cs_if_exist_p:N \patchcmd }
        { \cs_if_exist_p:N \__tblr_build_long_table:n }
        { \cs_if_exist_p:N \__tblr_halign_whole:Nn }
      }
      {
        \__docxlike_table_patch:Nnn \__tblr_build_long_table:n
          { \mode_leave_vertical: }
          { \docxlike_format_blank_para: \mode_leave_vertical: }
        \__docxlike_table_patch:Nnn \__tblr_halign_whole:Nn
          { \noindent }
          { \docxlike_format_blank_para: \noindent }
        \__docxlike_table_patch:Nnn \__tblr_build_long_table:n
          { \nobreak \noindent }
          { \nobreak \docxlike_format_blank_para: \noindent }
        \__docxlike_table_patch:Nnn \__tblr_build_long_table:n
          { \nobreak \noindent }
          { \nobreak \docxlike_format_blank_para: \noindent }
      }
  }
%    \end{macrocode}
%
% 单竖线计宽为零并居中于单元格边界，双线保持原行为；内部接口不存在时跳过补丁。
%    \begin{macrocode}
\cs_new_protected:Npn \__docxlike_table_tblr_vlines:
  {
    \bool_lazy_and:nnT
      { \cs_if_exist_p:N \__tblr_measure_and_update_vline_size:nn }
      { \cs_if_exist_p:N \__tblr_build_vline_segment_real:nn }
      {
        \cs_gset_eq:NN \__docxlike_table_tblr_vsize:nn \__tblr_measure_and_update_vline_size:nn
        \cs_gset_eq:NN \__docxlike_table_tblr_vseg:nn \__tblr_build_vline_segment_real:nn
        \cs_gset_protected:Npn \__tblr_measure_and_update_vline_size:nn ##1##2
          {
            \__docxlike_table_tblr_vsize:nn {##1} {##2}
            \bool_lazy_and:nnT
              { \l_docxlike_table_cell_bool }
              { \int_compare_p:nNn { \__tblr_spec_item:ne { vline } { [##2] / @vline-count } } = { 1 } }
              { \__tblr_spec_gput:nee { vline } { [##2] / @vline-width } { 0pt } }
          }
        \cs_gset_protected:Npn \__tblr_build_vline_segment_real:nn ##1##2
          {
            \bool_lazy_and:nnTF
              { \l_docxlike_table_cell_bool }
              { \int_compare_p:nNn { \__tblr_spec_item:ne { vline } { [##2] / @vline-count } } = { 1 } }
              {
                \hbox_set:Nn \l__docxlike_table_vline_box { \__docxlike_table_tblr_vseg:nn {##1} {##2} }
                \skip_horizontal:n { - 0.5 \box_wd:N \l__docxlike_table_vline_box }
                \box_use:N \l__docxlike_table_vline_box
                \skip_horizontal:n { - 0.5 \box_wd:N \l__docxlike_table_vline_box }
              }
              { \__docxlike_table_tblr_vseg:nn {##1} {##2} }
          }
      }
  }
\box_new:N \l__docxlike_table_vline_box
%</docxlike-table>
%    \end{macrocode}
%
% \Finale
